\section{Method and Decision Rules}
\label{sec:method}

We study a frozen model $M$ and a fixed action set $\mathcal{A}$. An input $x$
is rendered by a task adapter into a state string and a typed question; the
model returns a schema-validated decision with a probability. The policy reads
only the first call and decides whether to issue a second call. Model weights,
question templates and candidate label sets are never modified during search.
This isolates inference-time control from fine-tuning the backbone on the
routing objective.

\begin{figure*}[t]
\centering
\begin{tikzpicture}[
  font=\small,
  >=Stealth,
  node distance=6mm,
  box/.style={draw, rounded corners=2pt, align=center, inner sep=4pt,
              minimum height=8.5mm, fill=white},
  arr/.style={->, thick},
  dashedarr/.style={->, thick, dashed}
]
\node[box] (x) {input $x$};
\node[box, right=of x] (m) {frozen model $M$};
\node[box, right=of m] (p) {\noul{} call\\$p(x)$};
\node[box, right=of p] (g) {band gate\\$\ell \le p \le h$};
\node[box, right=of g, yshift=12mm] (c) {\choice{} call\\second view};
\node[box, right=of g, yshift=-12mm] (n) {\noul{} label\\no second call};
\draw[arr] (x) -- (m);
\draw[arr] (m) -- (p);
\draw[arr] (p) -- (g);
\draw[arr] (g) -- node[above, font=\scriptsize]{$u=1$} (c);
\draw[arr] (g) -- node[below left, font=\scriptsize]{else} (n);
\node[box, below=10mm of p, fill=gray!8] (fit) {$J(g)=Q(g)-\lambda\,C(g)$};
\node[box, left=of fit, fill=gray!8] (cache) {immutable output cache};
\node[box, right=of fit, fill=gray!8] (search) {GA / ES / NSGA-II};
\draw[dashedarr] (cache) -- (fit);
\draw[dashedarr] (fit) -- (search);
\end{tikzpicture}
\caption{The routing setting. A frozen model answers a first \noul{} question;
a two-parameter gate decides whether the same model is asked a second
\choice{} question. Search operates entirely on cached outputs: a candidate
genome is scored by replaying the cache, and the optimiser only sees the scalar
fitness $J$. The dashed path is offline; the solid path is the deployed policy.}
\label{fig:pipeline}
\end{figure*}

For the binary tasks, CodeXGLUE and FinDVer, the first action is \noul{}: the
model answers a yes/no question and returns $p(x)=P(\text{true}\mid x)$. The
second action is \choice{}: the model selects between two named labels. In
CodeXGLUE the questions are whether a function contains a vulnerability and
whether it should be labelled \texttt{secure} or \texttt{vulnerable}; in
FinDVer they are whether a claim is supported and whether it should be labelled
\texttt{supported} or \texttt{not supported}. The cascade always pays for
\noul{} first and pays for \choice{} only when the gate opens. A policy is a
genome $g_{\mathrm{cas}}=(\ell,h,u)$ with $\ell,h\in[0,1]$ and
$u\in\{0,1\}$:
\begin{equation}
\pi_{g}(x)=
\begin{cases}
\choice(x), & u=1\!\land\!\ell\le p(x)\le h,\\
\noul(x), & \text{otherwise}.
\end{cases}
\label{eq:cascade}
\end{equation}
The gate uses the positive-class probability rather than symmetric confidence,
because a \noul{} question has a designated true label. Genome repair clips
$\ell$ and $h$ to $[0,1]$, swaps them when $\ell>h$, and maps $u$ to $\{0,1\}$.
Two fixed policies are special cases: $u=0$ is \noul-only, and
$u=1,\ell=0,h=1$ is fixed switch-only, which still pays for the first \noul{}
call and therefore costs two calls per row.

ProcessBench has a different action space. The model returns a ranked list of
candidate steps, and the policy decides whether to answer with the top-1 step or
to abstain. A genome $g_{\mathrm{abs}}=(t,v,m)$ with $t,m\in[0,1]$ and
$v\in\{0,1\}$ uses the top-1 probability $s_1(x)$ and the top-1/top-2 margin
$\Delta(x)$:
\begin{equation}
\mathrm{abstain}(x)=
\begin{cases}
\text{true}, & s_1(x)<t,\\
\text{true}, & v=1\,\land\,\Delta(x)<m,\\
\text{false}, & \text{otherwise}.
\end{cases}
\label{eq:abstain}
\end{equation}
When the required confidence fields are missing, the evaluator fails closed:
the row is forced to answer and is counted as incorrect rather than silently
dropped. We report coverage (the fraction of rows answered) and selective
accuracy (accuracy on answered rows).

Fitness combines task quality with the cost of the second call:
\begin{equation}
J(g)=Q(g)-\lambda\,C(g),
\label{eq:fitness}
\end{equation}
where $Q$ is macro-F1 on the binary tasks and step accuracy on ProcessBench,
$C(g)$ is the mean number of model calls per row minus one, and $\lambda=0.05$
is fixed before search. Macro-F1 is primary because both binary datasets are
class-imbalanced; accuracy is reported alongside it. The first \noul{} call is
charged to every policy, including \noul-only, so the cost term isolates the
incremental call rather than the total inference bill. Model outputs are
generated once and replayed from an immutable JSON cache, which makes candidate
evaluation deterministic and cheap.

The GA uses tournament selection, uniform crossover, Gaussian mutation and bit
flips on $u$. The ES is a $(\mu+\lambda)$ strategy that mutates the two
endpoints and fixes $u=1$. NSGA-II maintains a non-dominated front, optimising
accuracy and negative call cost on the binary tasks and selective accuracy and
coverage on ProcessBench; because its quality objective is accuracy rather than
macro-F1, we keep its intervals separate. The main budget is 24 unique genome
evaluations with seeds 17, 23 and 41; a five-seed extension adds 59 and 73. A
probe raises the CodeXGLUE/Vela budget to 48. Random search uses the same budget
and cache.

One property of the policy class shapes the calibration experiment. The gate
depends on $p(x)$ only through the interval $\ell\le p\le h$. Any strictly
monotone transformation of $p$---temperature scaling, for example---maps the
interval to another interval and induces the same set of switched rows up to a
reparameterisation of the endpoints. A monotone recalibration therefore cannot
change the best achievable band policy on a fixed pool. It can move rows between
bands, but it cannot create a distinction between repairable and harmful rows
that the raw ordering did not already contain. Calibration and routing are
therefore evaluated separately.
